\documentclass[conference]{IEEEtran}
\usepackage[T1]{fontenc}
\usepackage[utf8]{inputenc}
\usepackage{amsmath,amssymb}
\usepackage{array}
\usepackage{booktabs}
\usepackage{multirow}
\usepackage{url}
\usepackage{microtype}
\newcolumntype{P}[1]{>{\raggedright\arraybackslash}p{#1}}
\setlength{\emergencystretch}{2em}
\title{Can We Stop Global Climate Change?\\A Data-to-Decision Framework for Deep Mitigation and Climate Risk Management}
\author{Research Plan Author Agent\\Xcientist four-agent simulation}
\begin{document}
\maketitle

\begin{abstract}
Stopping global climate change is not a single binary intervention. A scientifically useful formulation asks whether coordinated mitigation can drive net greenhouse-gas emissions toward zero while keeping residual physical risk and adaptation burden within declared limits under uncertainty. This report proposes the Climate Observing and Flexible Energy Transition framework (CO-FET), a coupled data-to-decision design that links global climate observations, probabilistic risk models, renewable-rich power-system planning, storage, demand flexibility, and governance triggers. The Survey stage establishes two bottlenecks: incomplete and unevenly shared observations, and fossil-fuel dependence combined with grid and storage constraints. The Idea stage selects CO-FET over an observation-only, technology-only, or carbon-removal-only route. The ExperimentDesign stage specifies a multi-region scenario experiment with held-out climate regimes, robust optimization, counterfactual energy portfolios, and explicit equity and reliability constraints. Mathematical formulations separate emissions trajectories, climate hazards, grid adequacy, and decision regret. The proposal makes a direct claim: global warming cannot be stopped by prediction alone, but a coordinated portfolio can sharply reduce additional warming and risk if observations are converted into enforceable mitigation and flexibility decisions. No empirical climate or deployment result is claimed; all outcomes are prospective and the study remains design-only.
\end{abstract}

\begin{IEEEkeywords}
climate change mitigation, Earth observation, renewable grids, energy storage, robust optimization, climate risk, decision support
\end{IEEEkeywords}

\section{Introduction}
The question of whether global climate change can be stopped is urgent but scientifically under-specified. Climate change is a coupled physical, ecological, technological, economic, and governance process. Stopping it cannot mean returning the climate system to a previous state on demand. A testable target is to reduce anthropogenic greenhouse-gas emissions to net zero, limit additional warming and hazards, and manage residual risk without transferring unacceptable burdens across regions or generations.

This report therefore reframes the question as follows: can a coordinated data-to-decision system identify robust mitigation portfolios that drive emissions toward net zero, preserve reliable and affordable energy, and reduce climate risk across plausible climate and socioeconomic futures? The proposed answer is the Climate Observing and Flexible Energy Transition framework (CO-FET).

The framing has two linked parts. First, the climate observing system must provide sustained, accessible, and quality-controlled information about the atmosphere, land, ocean, cryosphere, and human drivers. The Global Climate Observing System (GCOS), co-sponsored by WMO, IOC-UNESCO, UNEP, and ISC, defines Essential Climate Variables and states that observations support monitoring, understanding, prediction, climate services, and adaptation~\cite{gcos}. Observation is therefore an enabling infrastructure, not a substitute for mitigation.

Second, the energy system must move away from fossil fuels while maintaining adequacy. The International Energy Agency's Net Zero by 2050 roadmap describes a narrow pathway requiring large deployment of renewables, efficiency, electrification, innovation, and international cooperation~\cite{iea2021}. It also emphasizes that electricity becomes central and requires flexibility from batteries, demand response, hydrogen-based fuels, hydropower, and other resources. The IEA's grid assessment warns that grids can become a bottleneck when renewable projects and electrification expand faster than infrastructure, planning, and management~\cite{ieagrid}.

The key design claim is direct: data and low-carbon technologies become climate action only when joined by a robust decision policy. CO-FET couples observation quality with energy-system flexibility and climate-risk triggers. It tests portfolios against heat, water, food, coastal, infrastructure, and ecosystem risks rather than optimizing carbon alone. The goal is not to forecast one perfect future. It is to find actions that remain useful across a set of futures and reveal when a strategy is failing.

The report follows the four-agent workflow. The Survey Agent maps evidence and gaps. The Idea Agent selects CO-FET from competing routes. The ExperimentDesign Agent defines scenarios, equations, constraints, and falsification. The Author Agent freezes the proposal into an IEEE-style report. There are no observed deployment results, invented reductions, or claims that climate change has already been stopped.

\section{Evidence-Constrained Survey}
\subsection{The observing bottleneck}
GCOS states that a global climate observing system must make accurate and sustained observations available for monitoring, understanding, and predicting climate, and that Essential Climate Variables provide a systematic basis for observing change~\cite{gcos}. This supports a data layer with common variable definitions, uncertainty metadata, continuity, open access, and quality control. It also identifies the central limitation: an observation is useful only when it is timely, interoperable, and connected to a decision or service.

The survey therefore distinguishes coverage from usability. A dense network with inconsistent calibration may be less useful than a smaller network with traceable measurements. Satellite, in situ, ocean, land, and socioeconomic data have different spatial and temporal supports. Downstream models must retain these differences rather than treating every pixel or station as equally certain.

The data gap is not only technical. Countries and sectors differ in capacity to collect, curate, and share information. A global mitigation policy that uses data inaccessible to vulnerable regions can optimize the wrong risks. CO-FET makes data access and uncertainty part of the experiment rather than an afterthought.

\subsection{The mitigation and flexibility bottleneck}
The IEA Net Zero by 2050 roadmap describes the energy sector as central to the climate challenge and calls for a rapid expansion of clean energy, efficiency, electrification, innovation, and international cooperation~\cite{iea2021}. Its pathway is technologically diverse: renewables, electric transport, heat pumps, efficiency, sustainable bioenergy, hydrogen, carbon capture, and behavioral change have different timing, costs, land requirements, and reliability implications.

The IEA grid report adds an engineering constraint. Electricity demand expands as transport, buildings, and industry electrify, while renewable deployment requires more lines and better grid planning. The report describes delayed grid investment as a risk to clean-energy transition and electricity security~\cite{ieagrid}. This supports a design that evaluates generation, transmission, storage, demand response, and curtailment together. It rejects a simple substitution model in which variable renewable capacity is added without a reliability budget.

\subsection{Physical risk and adaptation}
The IPCC Sixth Assessment Synthesis Report provides the broad evidence base for climate impacts, adaptation, mitigation, and the interactions between them~\cite{ipccsyr}. Its synthesis supports a coupled risk formulation: hazard, exposure, and vulnerability must be considered together. The IPCC Working Group II report assesses impacts on ecosystems, biodiversity, human communities, infrastructure, health, and regional systems, as well as capacities and limits to adaptation~\cite{ipccwg2}.

This means that a portfolio can be low-carbon but still increase risk if it relies on water-intensive generation in a drought-prone basin, damages ecosystems, or shifts energy costs to households with little ability to adapt. CO-FET includes residual risk and distributional burden as constraints. It does not assume adaptation can compensate for unlimited warming.

\begin{table*}[t]
\centering
\caption{Survey evidence anchors and the claims permitted in the CO-FET design.}
\label{tab:evidence}
\small
\begin{tabular}{P{2.8cm}P{4.0cm}P{4.7cm}P{3.3cm}}
\toprule
Evidence anchor & Directly supported proposition & Design use & Forbidden inference \\
\midrule
GCOS~\cite{gcos} & Sustained, quality-controlled Essential Climate Variables support monitoring, prediction, services, and adaptation & Observation quality, access, uncertainty, and data continuity layer & More observations alone stop warming \\
IEA Net Zero 2050~\cite{iea2021} & A narrow pathway requires rapid clean-energy deployment, efficiency, electrification, innovation, and cooperation & Scenario portfolio and milestone constraints & One technology is sufficient \\
IEA Grids~\cite{ieagrid} & Grid upgrades, planning, and management are potential bottlenecks for renewable transition & Transmission, storage, demand response, and adequacy model & Installed renewable capacity equals delivered energy \\
IPCC AR6 Synthesis~\cite{ipccsyr} & Mitigation and adaptation must address interacting climate risks under uncertainty & Hazard, exposure, vulnerability, and regret metrics & Any single pathway is universally optimal \\
IPCC WGII~\cite{ipccwg2} & Risks affect ecosystems, communities, infrastructure, and regions with unequal adaptive capacity & Multi-sector risk and equity constraints & Adaptation removes all residual risk \\
\bottomrule
\end{tabular}
\end{table*}

\subsection{Research gap ledger}
Five gaps are accepted. First, observation quality is rarely evaluated by the decisions it changes. Second, energy-transition scenarios often treat grid flexibility as an implementation detail rather than a coupled resource. Third, climate risk is spatially and socially heterogeneous, so a global emissions objective can hide local harm. Fourth, policy decisions must be made before uncertainty is resolved; optimizing one forecast can create fragile infrastructure. Fifth, mitigation, adaptation, and data governance involve actors across countries, sectors, and agencies, but the decision interfaces between them are under-specified.

The survey rejects three overstrong claims. It rejects the claim that a better climate model alone can stop climate change. It rejects the claim that renewable capacity can be scaled without storage, transmission, demand flexibility, and reliability accounting. It rejects the claim that a global net-zero average guarantees equitable risk reduction. These rejections define the design space for a useful experiment.

\section{Idea Synthesis}
\subsection{Primary direction: CO-FET}
CO-FET is a coupled architecture with four layers. The observation layer ingests climate variables, energy-system measurements, land and water indicators, and data-quality metadata. The risk layer estimates hazards, exposure, vulnerability, and uncertainty for multiple sectors and regions. The energy layer simulates generation, transmission, storage, demand response, and fossil-fuel displacement. The policy layer selects robust milestones, investment triggers, and corrective actions.

The central hypothesis is:
\begin{quote}
A coordinated policy that jointly values emissions reduction, climate-risk reduction, grid reliability, affordability, equity, and observation quality will dominate carbon-only or technology-only policies in worst-case regret across plausible climate and socioeconomic futures.
\end{quote}

The claim is not that one optimization algorithm will solve governance. It is that a formal interface can make tradeoffs visible, test robust portfolios, and define when political and engineering commitments must be revisited.

\subsection{Competing routes}
An observation-first route expands monitoring and improves climate forecasts but leaves the fossil-fuel and grid decision problem under-specified. It is a necessary component, not a complete solution. A technology-only route selects the cheapest renewable or storage mix without explicitly modeling data uncertainty, distributional burden, or climate hazard. It can be efficient in one scenario and brittle elsewhere.

A carbon-removal-first route treats future negative emissions as permission to delay near-term mitigation. It is rejected as the primary direction because removal capacity, permanence, land, water, and monitoring create their own uncertainties. A single global carbon budget route is scientifically necessary but operationally incomplete if it does not specify regional risk, energy access, and reliability.

CO-FET combines the strongest elements: global observations, probabilistic risk, technology portfolios, flexibility, and policy triggers. It is deliberately not a claim that all sectors can be reduced to one score. The framework reports Pareto fronts and robust actions rather than a falsely precise ranking of countries.

\subsection{Predictions}
CO-FET predicts that robust portfolios will include a mix of renewable generation, transmission, storage, demand response, efficiency, and selected firm low-carbon resources rather than one dominant technology. It predicts that the value of observation increases when it reduces decision regret or detects threshold crossing. It predicts that portfolios optimized only for mean emissions will underperform on reliability or regional risk in some futures. It predicts that the best policy is staged: commit to no-regret actions early, preserve options, and trigger stronger measures when observed evidence crosses predeclared thresholds.

The framework also predicts a meaningful tradeoff. Climate data infrastructure and grid flexibility require investment before their full value is visible. The experiment therefore evaluates option value, not only immediate cost. A portfolio that costs slightly more in the median scenario may be superior if it avoids catastrophic shortfalls across plausible futures.

\section{ExperimentDesign}
\subsection{Study object and scenario ensemble}
The study object is a policy portfolio, not a forecast. A portfolio contains emissions measures, generation and storage investments, transmission expansion, demand response, observation upgrades, and risk-management actions. It is evaluated across an ensemble of climate, socioeconomic, technology, and governance scenarios. Each scenario preserves physical uncertainty, demand uncertainty, renewable availability, cost uncertainty, and data quality.

The spatial unit is a set of connected regions $r$ with energy nodes, climate-risk zones, and exposed populations. The time unit is a planning interval $t$ with finer operational substeps for electricity dispatch. The horizon includes near-term milestones and long-term outcomes. Decisions are staged: irreversible investments are distinguished from flexible options, and the policy can update after an observation or trigger.

The scenario ensemble has four layers. Climate trajectories provide temperature, precipitation, extremes, sea level, and compound hazards. Socioeconomic trajectories provide demand, population, land use, and vulnerability. Energy trajectories provide technology costs, outages, resource availability, storage performance, and transmission limits. Observation trajectories provide missingness, calibration error, spatial coverage, and reporting delay. A scenario is not treated as a prediction; it is a stress test for policy robustness.

\subsection{Emissions accounting}
Let $E_{r,t,k}$ denote activity in region $r$, period $t$, and source or technology $k$, and let $g_k$ be its life-cycle or direct emissions factor under a declared accounting boundary. Net emissions are:
\begin{equation}
G_t=\sum_r\sum_k g_k E_{r,t,k}-R_t,
\label{eq:emissions}
\end{equation}
where $R_t$ is verified removal within the scenario. The primary mitigation objective is to reduce $G_t$ toward zero while reporting gross emissions, removal dependence, and cumulative emissions separately. This prevents a portfolio from appearing successful solely because uncertain future removal is credited in advance.

The emissions layer includes energy demand, industrial process emissions, land-use change, methane and other gases through a common forcing-equivalent convention declared before analysis. Sensitivity runs vary the accounting boundary. A portfolio is not accepted because it reaches a nominal annual target if cumulative emissions, leakage, or rebound effects violate the declared boundary.

\subsection{Climate risk model}
For region $r$, sector $s$, time $t$, and scenario $w$, risk is represented as:
\begin{equation}
Q_{r,s,t,w}=H_{r,s,t,w}\cdot X_{r,s,t,w}\cdot V_{r,s,t,w}\cdot (1-A_{r,s,t,w}),
\label{eq:risk}
\end{equation}
where $H$ is hazard intensity or probability, $X$ is exposure, $V$ is vulnerability, and $A$ is effective adaptation capacity on a normalized scale. The product is a transparent index, not a claim that all hazards obey a linear relationship. Sector-specific damage functions are tested as sensitivity models.

The risk vector includes heat, water scarcity or flood, food, coastal, infrastructure, ecosystem, and health-relevant indicators where data quality permits. The primary risk statistic is a weighted distribution of residual loss, with weights declared by region and population rather than chosen after seeing results. Distributional metrics report the worst affected region, the lower access group, and the difference between high- and low-capacity regions.

\subsection{Power-system balance and flexibility}
For each node and operational substep, electricity balance requires:
\begin{equation}
\sum_k P_{r,t,k}+D^{\mathrm{imp}}_{r,t}-D^{\mathrm{exp}}_{r,t}
=D_{r,t}+\Delta S_{r,t}+L_{r,t},
\label{eq:balance}
\end{equation}
where $P$ is generation or discharge, $D^{\mathrm{imp}}$ and $D^{\mathrm{exp}}$ are imports and exports, $D$ is demand, $\Delta S$ is net storage charging, and $L$ is modeled loss. Resource availability, ramping, minimum stable output, transmission capacity, storage state of charge, and reserve requirements are explicit constraints.

Reliability is reported through unserved energy, loss-of-load probability, reserve margin, curtailment, and recovery time after an extreme event. Demand response is a bounded shift with rebound and comfort constraints, not free energy. Storage is evaluated for duration, degradation, round-trip efficiency, material availability, and replacement timing. Hydrogen and firm low-carbon resources are included as scenario options rather than assumed solutions.

\subsection{Robust optimization}
Let $z$ be a staged policy and $J_w(z)$ the vector of cost, emissions, reliability loss, climate risk, affordability burden, and data-access burden in scenario $w$. A minimax-regret objective is:
\begin{equation}
\begin{aligned}
z^*=\mathop{\mathrm{arg\,min}}\limits_{z\in\mathcal{Z}}\; &\max_{w\in\mathcal{W}}\nonumber\\
&\left\|J_w(z)-\min_{z'\in\mathcal{Z}}J_w(z')\right\|_\infty\\
\text{subject to }&G_t\rightarrow 0,\nonumber\\
&\operatorname{Reliability}(z,w)\geq \underline{q},\nonumber\\
&\operatorname{Burden}_{r}(z,w)\leq \overline{b}_r.
\end{aligned}
\label{eq:robust}
\end{equation}
Here $\mathcal{Z}$ is the feasible portfolio set, $\mathcal{W}$ is the scenario set, and $\underline{q}$ and $\overline{b}_r$ are predeclared reliability and burden limits. The vector objective is solved with Pareto analysis and documented weights. Robustness means low regret across scenarios, not optimality in a single forecast.

\begin{sloppypar}
The baseline policies are carbon-only optimization, technology-only cost minimization, observation-first investment, and delayed action.

CO-FET is compared with the same physical constraints and scenario ensemble. A policy is not called robust if it relies on an untested assumption about future technology, global coordination, or removal permanence.
\end{sloppypar}

\subsection{Value of observation}
Observation investment $o$ changes uncertainty and delay. Its decision value is defined by the reduction in expected or worst-case regret:
\begin{equation}
\operatorname{VOI}(o)=
\mathbb{E}_{y\sim p(y|o)}
\left[\min_z \mathbb{E}_{w|y}J_w(z)\right]
-\min_z \mathbb{E}_{w}J_w(z),
\label{eq:voi}
\end{equation}
with sign convention reported so that positive value denotes an improvement after paying observation cost. The value is also evaluated for threshold decisions: whether an observation changes a grid investment, a risk trigger, or a mitigation milestone.

This formulation makes the observing system accountable to decisions. A variable can be scientifically interesting but have low policy value if it does not reduce uncertainty where action changes. Conversely, a modest local measurement can have high value if it prevents a reliability or exposure mistake. Coverage, latency, calibration, openness, and equity are part of $o$.

\subsection{Adaptive triggers}
The policy is staged using a trigger map. At review time $t$, the system compares observed indicators with a predeclared envelope. If emissions, reliability, hazard, or vulnerability indicators cross a threshold, the policy escalates a previously approved action. If observations remain within the envelope, the policy continues while preserving option value. A trigger cannot authorize a policy outside legal or democratic governance; it only defines when a review is mandatory.

Each trigger records indicator definition, baseline, uncertainty interval, response owner, lead time, and rollback option. Trigger hysteresis prevents repeated switching around a noisy threshold. The analysis compares fixed, reactive, and robust-staged policies. The desired result is earlier and less costly correction, not maximal intervention.

\subsection{Evaluation metrics}
Primary metrics are worst-case regret, cumulative net emissions, residual multi-sector risk, reliability loss, affordability burden, and access to observation data. Secondary metrics include cost per avoided unit of emissions, renewable curtailment, storage utilization, transmission congestion, adaptation co-benefits, decision latency, and option value.

The analysis reports Pareto fronts and scenario envelopes. It does not compress all values into one undisclosed score. Weight sensitivity, region-weight sensitivity, removal accounting, climate sensitivity, and technology availability are required. A policy that wins only for one weight vector is labeled preference-dependent.

\subsection{Controls and ablations}
The experiment includes four ablations. Removing the observation layer tests whether data have decision value. Removing grid flexibility tests whether capacity-only decarbonization is brittle. Removing risk and equity constraints tests whether carbon-only optimization hides burden. Removing adaptive triggers tests whether staging improves regret.

Two negative controls protect interpretation. First, a perfect-forecast policy estimates the upper bound of information value but is not achievable. Second, a naive capacity-substitution policy adds renewable nameplate capacity without storage or grid reinforcement. If CO-FET cannot outperform these controls under their declared assumptions, its claim is weakened.

\subsection{Analysis sequence}
The analysis follows eight steps:
\begin{enumerate}
\item Freeze scenario definitions, data provenance, accounting boundaries, weights, thresholds, and reliability limits.
\item Construct climate, socioeconomic, energy, and observation scenario ensembles.
\item Calibrate baseline energy dispatch and climate-risk models against held-out periods.
\item Optimize carbon-only, technology-only, observation-first, delayed, and CO-FET policies.
\item Evaluate each policy on held-out scenarios, extreme events, and distributional outcomes.
\item Run ablations for observations, flexibility, risk constraints, and adaptive triggers.
\item Perform sensitivity, value-of-information, and model-discrepancy analyses.
\item Report Pareto fronts, worst-case regret, trigger behavior, and conditions for falsification.
\end{enumerate}

\section{Expected Interpretations and Falsification}
\subsection{Interpretation matrix}
The experiment is designed to produce a decision-relevant result rather than a technology leaderboard. Four outcomes are pre-specified.

A robust-transition outcome occurs if CO-FET reduces worst-case regret and residual risk while meeting reliability, affordability, and equity limits across held-out scenarios. This would support coupled observation, mitigation, flexibility, and staged governance.

A data-value outcome occurs if observation investment changes decisions and reduces regret, while energy-system policies without improved observations fail under uncertainty. This would show that climate data infrastructure earns its value through decisions, not through coverage alone.

A flexibility-dominant outcome occurs if storage, transmission, demand response, and firm low-carbon resources explain most of the advantage while observation upgrades add little. This would redirect investment toward power-system flexibility, without invalidating observation for physical risk management.

A carbon-only failure occurs if emissions optimization achieves a favorable mean but produces unacceptable reliability, regional risk, or distributional burden in held-out scenarios. Such a result would falsify the assumption that net emissions alone is an adequate objective.

\begin{table}[t]
\centering
\caption{Pre-registered interpretation rules for CO-FET.}
\label{tab:decisions}
\footnotesize
\begin{tabular}{P{2.0cm}P{2.5cm}P{2.5cm}}
\toprule
Outcome & Scenario signature & Interpretation \\
\midrule
Robust transition & Regret and risk fall; constraints hold & Coupled policy is supported \\
Data value & Observations change decisions and reduce regret & Data infrastructure earns operational value \\
Flexibility & Grid resources dominate incremental benefit & Reliability is the binding transition constraint \\
Carbon-only failure & Mean emissions improve but risk or reliability fails & Single-objective policy is rejected \\
\bottomrule
\end{tabular}
\end{table}

\subsection{Explicit falsification}
CO-FET is falsified if it cannot outperform a carbon-only or technology-only policy on worst-case regret after all policies receive the same scenario information and physical constraints. It is also falsified if its advantage disappears under held-out climate regimes, if its policy relies on unrealistically perfect observations, or if reliability and affordability constraints cannot be met without hidden exclusions.

The observation layer is not retained as a decision component if additional observations do not change any mitigation, grid, or risk-management action and do not reduce uncertainty in a decision-relevant region. The flexibility layer is not retained if a capacity-only policy meets reliability under all declared scenarios without additional storage, transmission, demand response, or firm resources. The adaptive trigger layer is not retained if it increases switching, regret, or delay relative to a fixed robust plan.

A negative result is not converted into a claim that climate science or renewable energy is ineffective. It identifies which link in the data-to-decision chain failed. The framework reports model discrepancy, scenario dependence, and distributional effects rather than hiding them in one global average.

\subsection{Equity and uncertainty}
A global reduction in emissions can coexist with local energy poverty or concentrated climate risk. CO-FET therefore reports cost, reliability, and residual risk by region and access group. The policy is constrained from meeting global targets by shifting unserved energy or adaptation burden to a low-capacity region. Weight sensitivity is mandatory because equity priorities are normative inputs, not parameters that can be discovered from data.

Uncertainty is an input to action. A broad scenario envelope should favor flexible, reversible, and no-regret investments, while a threshold crossing should trigger a predefined review. The policy does not treat uncertainty as a reason for inaction; it treats it as a reason to preserve options and monitor the variables that can change the decision.

\section{Implementation and Governance}
\subsection{Data-to-decision architecture}
The project begins with a common data and provenance layer. Each Essential Climate Variable, energy measurement, socioeconomic indicator, and risk estimate receives a definition, spatial and temporal support, quality flag, uncertainty interval, and update time. The data system records whether a variable is open, restricted, delayed, or missing. This makes the observation bottleneck measurable.

The second layer produces scenario inputs and risk indicators. The third layer solves dispatch and investment decisions under the same constraints for every competing policy. The final layer presents Pareto fronts, scenario envelopes, trigger status, and distributional results to decision makers. No model output is treated as a legal mandate. The policy layer identifies commitments and review points that must be assigned to accountable institutions.

\subsection{Staged policy}
CO-FET separates actions into three classes. Immediate actions are low-regret measures such as efficiency, data continuity, maintenance, and grid planning. Option-preserving actions include modular storage, transmission corridors, demand-response programs, and interoperable observation infrastructure. Triggered actions are larger investments or restrictions whose timing depends on emissions, reliability, hazard, or vulnerability indicators.

Every trigger records its indicator, uncertainty interval, threshold, lead time, response owner, and rollback option. Hysteresis prevents repeated activation around a noisy boundary. Trigger rules are evaluated in simulation before adoption. A trigger can require a public review; it cannot silently bypass national or local governance.

\subsection{Coordination protocol}
The scenario and portfolio registry records assumptions shared across climate scientists, energy planners, engineers, public agencies, and affected communities. A change to an emissions factor, reliability standard, hazard model, or burden weight creates a versioned event. Results are not compared across versions without preserving the corresponding assumptions.

The coordination experiment compares a shared-registry policy with an uncoordinated policy in which regions use inconsistent observations and planning horizons. The outcome is decision regret, not a score for institutional virtue. If coordination does not improve decisions under the tested assumptions, the model records the boundary rather than presuming cooperation.

\subsection{Sample and computational power}
This is a scenario and optimization experiment, so power is assessed as model-recovery and scenario-coverage power. The simulation varies the number of climate and socioeconomic trajectories, extreme-event frequency, technology cost uncertainty, renewable correlation, storage duration, transmission expansion, and observation error. A design is retained only if it can recover a known robust portfolio when one exists and reject a fragile policy when it fails in held-out scenarios.

The scenario ensemble is split before optimization into development, stress-test, and held-out sets. Hyperparameters, weights, and trigger thresholds are fixed within the development stage. Stress tests probe extremes; held-out scenarios estimate generalization. If the policy wins only on the scenarios used to construct it, it is not considered robust.

\subsection{Environmental and social safeguards}
A low-carbon portfolio can create land, water, mineral, biodiversity, or community burdens. These impacts are included as constraints or reported co-impacts when data permit. The design does not assume that a technology is sustainable simply because it has a low operational carbon factor. Any deployment proposal requires local consultation, permitting, environmental review, and distributional assessment outside the model.

\section{Risks and Reproducibility}
\subsection{Model discrepancy}
Climate models, energy models, and risk functions are imperfect. CO-FET uses multiple model structures and reports discrepancy between them. A policy is not labeled robust because it survives only one simulator. If model disagreement changes the recommended portfolio, the disagreement becomes a decision-relevant uncertainty and an observation priority.

\subsection{Accounting and leakage}
Net emissions can be distorted by system boundaries, imported goods, land-use assumptions, or uncertain removals. Gross emissions, cumulative emissions, removal dependence, and leakage are reported separately. Sensitivity analyses vary the accounting boundary and removal permanence. A portfolio cannot claim success by shifting emissions outside its modeled boundary.

\subsection{Distribution shift}
Renewable resource patterns, demand, technology costs, grid outages, and climate extremes can change. Held-out climate regimes and time slices test shift. The observation layer monitors drift in measurements and model residuals. If a variable leaves its calibration envelope, the policy moves to review rather than producing a false precise recommendation.

\subsection{Optimization and political risks}
A mathematically optimal portfolio can be infeasible because of permitting, finance, supply chains, or political coordination. Feasibility constraints include construction rates, material availability, lead times, and regional implementation capacity where data permit. The portfolio registry reports which assumptions are engineering, economic, or governance commitments. This prevents a scenario result from being misrepresented as a guaranteed implementation plan.

\subsection{Reproducibility}
The repository contains the scenario manifest, observation schema, emissions factors, risk functions, grid model, optimization code, policy baselines, ablation scripts, and random seeds. Synthetic scenarios are included so that the robust-optimization and value-of-observation calculations can be reproduced without restricted data. Every figure and table is generated from a versioned scenario set. The project reports null findings, failed triggers, infeasible portfolios, and distributional harms.

\subsection{Scope of this report}
This is a research proposal, not a climate projection or a policy commitment. It reports no measured emissions reduction, no observed grid reliability improvement, and no claim that a particular country or technology will solve climate change. Its strong conclusion is methodological and actionable: the most credible path to stopping additional anthropogenic warming is a coordinated portfolio that combines rapid mitigation with grid flexibility and a decision-linked observing system.

\section{Conclusion}
Global climate change cannot be stopped by better prediction alone, and it cannot be stopped by adding renewable nameplate capacity without system flexibility. A scientifically meaningful target is to stop the growth of anthropogenic forcing by driving net greenhouse-gas emissions toward zero while reducing residual risk and preventing mitigation burdens from being shifted to vulnerable regions.

CO-FET makes this target testable. It links sustained climate observations to uncertainty-aware hazard and vulnerability estimates, couples renewable deployment to storage, transmission, demand response, and reliability, and uses robust optimization to preserve useful actions across plausible futures. Its trigger layer converts new observations into staged reviews rather than treating every forecast as a command.

The expected result is a portfolio rather than a miracle technology: early efficiency and clean-energy deployment, accelerated grid modernization, diverse flexibility resources, protected observation infrastructure, and risk-aware milestones. In some regions, the binding constraint will be transmission or storage. In others, it will be access, finance, water, land, or vulnerability. The framework is designed to expose those differences.

The answer to the original question is therefore conditional but direct. We can stop additional human-driven climate change in the operational sense of reaching net-zero emissions, but only if scientific information is converted into coordinated, enforceable, and adaptable decisions. CO-FET should be accepted only if it demonstrates lower worst-case regret and residual risk than simpler policies without violating reliability, affordability, or equity constraints. If it fails that test, the failed layer is more informative than a slogan.

\begin{thebibliography}{00}
\bibitem{gcos} Global Climate Observing System, Global Climate Observing System, World Meteorological Organization, 2026. [Online]. Available: \url{https://gcos.wmo.int/}.
\bibitem{ipccsyr} Intergovernmental Panel on Climate Change, Climate Change 2023: Synthesis Report, IPCC, Geneva, Switzerland, 2023. [Online]. Available: \url{https://www.ipcc.ch/report/ar6/syr/}.
\bibitem{ipccwg2} Intergovernmental Panel on Climate Change, Climate Change 2022: Impacts, Adaptation and Vulnerability, IPCC, Cambridge, U.K., 2022. [Online]. Available: \url{https://www.ipcc.ch/report/ar6/wg2/}.
\bibitem{iea2021} International Energy Agency, Net Zero by 2050: A Roadmap for the Global Energy Sector, IEA, Paris, France, 2021. [Online]. Available: \url{https://www.iea.org/reports/net-zero-by-2050}.
\bibitem{ieagrid} International Energy Agency, Electricity Grids and Secure Energy Transitions, IEA, Paris, France, 2023. [Online]. Available: \url{https://www.iea.org/reports/electricity-grids-and-secure-energy-transitions}.
\end{thebibliography}
\end{document}
